% relazione.tex — Client MCP + orchestratore Langgraph per l'advisor di
% mobilità Snap4City (UNIFI, Sistemi Distribuiti, elaborato Tipo A).
%
% Pacchetti richiesti: inputenc, fontenc, babel(italian), geometry, listings,
% xcolor, booktabs, tikz(librerie arrows.meta, positioning), hyperref.
% Compilazione:  pdflatex relazione.tex   (oppure: latexmk -pdf relazione.tex)
% Se TikZ non è disponibile, sostituire l'ambiente tikzpicture della Fig. 1 con
% il blocco {verbatim} di fallback indicato nei commenti vicino alla figura.
\documentclass[11pt,a4paper]{article}

\usepackage[utf8]{inputenc}
\usepackage[T1]{fontenc}
\usepackage[italian]{babel}
\usepackage[margin=2.5cm]{geometry}
\usepackage{xcolor}
\usepackage{booktabs}
\usepackage{listings}
\usepackage{tikz}
\usetikzlibrary{arrows.meta, positioning}
\usepackage[hidelinks]{hyperref}

% --- stile per i frammenti di codice -----------------------------------------
\definecolor{codebg}{rgb}{0.97,0.97,0.97}
\definecolor{codecomment}{rgb}{0.40,0.55,0.40}
\definecolor{codekw}{rgb}{0.10,0.30,0.65}
\definecolor{codestr}{rgb}{0.60,0.30,0.10}
\lstdefinestyle{py}{
  language=Python,
  backgroundcolor=\color{codebg},
  basicstyle=\ttfamily\footnotesize,
  keywordstyle=\color{codekw},
  commentstyle=\color{codecomment},
  stringstyle=\color{codestr},
  numbers=left, numberstyle=\tiny\color{gray}, numbersep=6pt,
  breaklines=true, frame=single, framesep=4pt, showstringspaces=false,
  columns=fullflexible, keepspaces=true,
}
\lstset{style=py}

\title{\textbf{Client MCP e orchestratore Langgraph\\ per l'advisor di mobilità Snap4City}}
\author{Elaborato Tipo A --- Sistemi Distribuiti\\ Università degli Studi di Firenze}
\date{\today}

\begin{document}
\maketitle

\begin{abstract}
\noindent
Questo elaborato descrive un \emph{client} per il servizio di consulenza alla
mobilità di Snap4City. L'utente pone in linguaggio naturale una domanda di
spostamento o di trasporto pubblico; un grafo \emph{deterministico} realizzato
con Langgraph guida il flusso \texttt{understand $\rightarrow$ execute
$\rightarrow$ respond} e restituisce un oggetto JSON (widget) pronto per essere
reso da una dashboard Snap4City. Il modello linguistico Llama4 interviene solo
in due punti --- l'estrazione strutturata dei parametri della richiesta e la
formulazione testuale della risposta --- mentre l'orchestrazione degli strumenti
(geocodifica, calcolo del percorso, scoperta del trasporto pubblico) è guidata
da codice Python. Il vero server MCP che implementa gli strumenti è gestito dal
referente ed è raggiungibile sulla rete interna; l'elaborato consegna quindi il
\emph{client}, l'orchestratore e l'interfaccia di test. Si documentano
l'architettura, le scelte progettuali con le relative \emph{motivazioni}, e le
limitazioni note lato server emerse dalla sperimentazione.
\end{abstract}

\tableofcontents
\bigskip

% =============================================================================
\section{Introduzione}
% =============================================================================
Snap4City è una piattaforma per smart city che espone, fra l'altro, servizi di
mobilità sull'area della Toscana (geocodifica, routing punto-punto, dati del
trasporto pubblico locale). L'obiettivo del lavoro è offrire un'interfaccia
conversazionale a questi servizi: l'utente scrive una frase come
``\textit{da Piazza del Duomo a Santa Croce a piedi}'' e ottiene una risposta in
linguaggio naturale, accompagnata da un payload strutturato (distanza, durata,
geometria del percorso in formato WKT) che la dashboard può visualizzare.

\paragraph{Confine della consegna.} Gli strumenti veri e propri risiedono sul
server MCP del referente, distribuito sulla rete interna di Snap4City
(accesso diretto da JupyterHub all'indirizzo \texttt{192.168.1.117:8000}). Il
presente elaborato consegna esclusivamente: (i) il \emph{client} MCP, (ii)
l'orchestratore Langgraph, (iii) una REPL da terminale (\texttt{chat.py}) come
collante di test. Nessuno strumento è implementato lato client.

\paragraph{Stack tecnologico.} La Tabella~\ref{tab:stack} riassume le tecnologie
impiegate.

\begin{table}[h]
\centering
\begin{tabular}{ll}
\toprule
\textbf{Componente} & \textbf{Tecnologia} \\
\midrule
Linguaggio          & Python 3.10+ (ambiente conda 3.11 su JupyterHub) \\
Client MCP          & FastMCP 2.x \\
Orchestrazione      & Langgraph 1.x (\texttt{StateGraph}) \\
Trasporto           & HTTP Streamable verso la dashboard \\
Modello linguistico & Llama4 (endpoint \texttt{llama4-agentic-inference}, OpenAI-compatibile) \\
Rendering           & widget della dashboard Snap4City (esterno al client) \\
\bottomrule
\end{tabular}
\caption{Stack tecnologico del client.}
\label{tab:stack}
\end{table}

% =============================================================================
\section{Architettura del sistema}
% =============================================================================
Il cuore del sistema è un grafo \emph{lineare e deterministico} composto da tre
nodi (Fig.~\ref{fig:graph}). Una sola richiesta utente attraversa il grafo dal
nodo di ingresso fino a \texttt{END}.

\begin{figure}[h]
\centering
% --- Se TikZ non è disponibile, sostituire questo tikzpicture con: ------------
% \begin{verbatim}
% [messages] -> understand -> execute -> respond -> END -> widget JSON
%               (LLM,slot)   (Python)    (LLM,prosa)
% \end{verbatim}
% -----------------------------------------------------------------------------
\begin{tikzpicture}[
  node distance=14mm and 16mm,
  box/.style={draw, rounded corners, minimum height=10mm, minimum width=22mm, align=center, font=\small},
  llm/.style={box, fill=blue!8},
  py/.style={box, fill=green!8},
  io/.style={box, fill=gray!10},
  arr/.style={-{Stealth[length=2.5mm]}, thick},
]
  \node[io]  (in)   {messages\\(storia)};
  \node[llm] (und)  [right=of in]  {understand\\\scriptsize LLM: slot};
  \node[py]  (exe)  [right=of und] {execute\\\scriptsize Python: tool};
  \node[llm] (res)  [right=of exe] {respond\\\scriptsize LLM: prosa};
  \node[io]  (out)  [right=of res] {widget\\JSON};
  \draw[arr] (in)  -- (und);
  \draw[arr] (und) -- (exe);
  \draw[arr] (exe) -- (res);
  \draw[arr] (res) -- (out);
\end{tikzpicture}
\caption{Il grafo deterministico \texttt{understand $\rightarrow$ execute
$\rightarrow$ respond}. In azzurro i nodi con LLM, in verde il nodo di pura
logica Python.}
\label{fig:graph}
\end{figure}

\paragraph{Motivazione progettuale: perché un grafo deterministico.}
La scelta architetturale fondante è che \emph{l'LLM non sceglie mai gli
strumenti}. Un primo prototipo usava il modello in modalità agentica
(\texttt{tool\_choice="auto"}): si è rivelato inaffidabile perché Llama4, quando
``vuole raccontare'' (ad esempio dopo un fallimento del routing), emette la
chiamata allo strumento come \emph{testo pythonico} (\texttt{[routing(\dots)]})
anziché come \texttt{tool\_calls} strutturati. Il testo trapelava nella risposta
finale, il campo dati restava vuoto e lo strumento non veniva realmente
eseguito. La soluzione definitiva è stata eliminare il ciclo agentico: il
modello sceglie soltanto i \emph{parametri} (in modo forzato e strutturato) e la
\emph{prosa} (senza strumenti), mentre Python guida deterministicamente le
chiamate. Questo modo di operare elimina alla radice la classe di errori
descritta.

% =============================================================================
\section{Lo stato condiviso (\texttt{AdvisorState})}
% =============================================================================
I nodi comunicano attraverso un unico oggetto di stato, un \texttt{TypedDict}
(Tabella~\ref{tab:state}).

\begin{table}[h]
\centering
\begin{tabular}{ll}
\toprule
\textbf{Campo} & \textbf{Significato} \\
\midrule
\texttt{messages}     & storia della conversazione (system/user/assistant), per il multi-turno \\
\texttt{intent}       & \texttt{route} $\mid$ \texttt{tpl\_lines/routes/stops/timeline} $\mid$ \texttt{other} \\
\texttt{slots}        & output di \texttt{understand} (interpretazione strutturata) \\
\texttt{tool\_results} & audit \texttt{[\{name, args, result\}]} di ogni chiamata \\
\texttt{unsupported}  & vero quando \texttt{execute} non può eseguire un flusso deterministico \\
\texttt{final}        & il widget JSON assemblato da \texttt{respond} \\
\bottomrule
\end{tabular}
\caption{Campi di \texttt{AdvisorState}.}
\label{tab:state}
\end{table}

\paragraph{Motivazione progettuale: dati derivati fuori dallo stato.} Langgraph
riconosce \emph{solo} le chiavi dichiarate nello schema: una chiave non
dichiarata restituita da un nodo viene scartata silenziosamente. Per questo i
dati derivati non vengono aggiunti allo stato; ad esempio gli ``slot mancanti''
non hanno un canale dedicato ma sono ricalcolati nel nodo \texttt{respond} a
partire dal canale \texttt{slots} già presente. Lo stato multi-turno, invece,
\emph{esiste già} nel campo \texttt{messages}: per proseguire una conversazione
è sufficiente ripassare i \texttt{messages} del turno precedente come storia.

% =============================================================================
\section{Nodo \texttt{understand} --- estrazione strutturata (LLM)}
% =============================================================================
Il primo nodo legge l'ultimo turno dell'utente e ne estrae i parametri tramite
una chiamata di funzione \emph{forzata}. La funzione \texttt{extract\_slots}
(Listato~\ref{lst:schema}) non è un vero strumento MCP: serve unicamente a
costringere il modello a produrre un output strutturato. La classificazione
avviene su due assi --- \texttt{request\_type} (journey / transit\_info / other)
e \texttt{info\_kind} (lines / routes / stops / timeline) --- poi ridotti a un
unico \texttt{intent} interno.

\begin{lstlisting}[caption={Lo schema (sintetico) della funzione che forza
l'output strutturato.}, label={lst:schema}]
resp = await llm.achat(
    messages=[{"role": "system", "content": UNDERSTAND_SYSTEM}, *convo],
    tools=[_EXTRACT_SLOTS_SCHEMA],
    tool_choice={"type": "function", "function": {"name": "extract_slots"}},
    temperature=0,
)
# campi: request_type, info_kind, origin_text, destination_text, mode,
#        agency_text, line_text, stop_text   (tutti "required")
\end{lstlisting}

\paragraph{Motivazione progettuale.}
\begin{itemize}
  \item \textbf{\texttt{temperature=0}}: l'estrazione dei parametri deve essere
  deterministica e riproducibile --- a parità di frase, gli stessi slot. Qui non
  serve alcuna creatività; ogni componente casuale sarebbe solo rumore.
  \item \textbf{\texttt{tool\_choice} forzato}: obbliga il gateway a restituire
  \texttt{tool\_calls} strutturati, eliminando il percorso del ``testo
  pythonico'' che affliggeva l'uso libero degli strumenti.
  \item \textbf{tutti i campi \texttt{required}}: Llama4 tende a riempire solo i
  parametri obbligatori, omettendo silenziosamente gli opzionali (in una prova
  reale estraeva l'origine ma perdeva la destinazione). Dichiarando tutti i
  campi obbligatori, una stringa vuota \texttt{''} segnala in modo controllato
  uno slot effettivamente assente.
  \item \textbf{estrazione del solo testo del luogo}: il modello non produce mai
  coordinate; la conversione in coordinate è compito di uno strumento separato.
\end{itemize}

% =============================================================================
\section{Nodo \texttt{execute} --- esecuzione deterministica (Python)}
% =============================================================================
Il secondo nodo non usa alcun LLM: esegue la catena di strumenti corrispondente
all'intento.

\paragraph{Flusso \texttt{route}.} Geocodifica dell'origine, geocodifica della
destinazione, quindi chiamata a \texttt{routing} con la modalità richiesta. Ogni
chiamata è registrata in \texttt{tool\_results}. Per le modalità a piedi, un
fallimento comporta un \emph{unico} nuovo tentativo con l'altro profilo pedonale
(i due profili percorrono percorsi diversi nel grafo, quindi uno può riuscire
dove l'altro restituisce un corpo vuoto). Il ripiego avviene una sola volta: il
profilo richiesto ha già esaurito la ``scala'' completa dei ritentativi per il
caso transitorio, quindi il transitorio è escluso e si prova soltanto l'altro
percorso del grafo. Per le modalità auto / trasporto pubblico, invece, non si
cambia modalità automaticamente: in caso di fallimento sarà il nodo
\texttt{respond} a suggerire un'alternativa, per non alterare l'intento
dell'utente.

\paragraph{Flussi \texttt{tpl\_*} (trasporto pubblico).} Sono catene di scoperta
(Tabella~\ref{tab:tpl}) delegate al modulo \texttt{tpl.py}.

\begin{table}[h]
\centering
\begin{tabular}{ll}
\toprule
\textbf{Intento} & \textbf{Catena} \\
\midrule
\texttt{tpl\_lines}    & agenzie $\rightarrow$ risoluzione agenzia $\rightarrow$ \texttt{tpl\_lines(agency)} \\
\texttt{tpl\_routes}   & \dots $\rightarrow$ \texttt{tpl\_routes\_by\_line(line, agency)} \\
\texttt{tpl\_stops}    & \dots $\rightarrow$ \texttt{tpl\_stops\_by\_route(route)} (prime 2 tratte) \\
\texttt{tpl\_timeline} & \dots $\rightarrow$ match del nome fermata $\rightarrow$ \texttt{tpl\_stop\_timeline(stop)} \\
\bottomrule
\end{tabular}
\caption{Le quattro catene di scoperta del trasporto pubblico locale.}
\label{tab:tpl}
\end{table}

\paragraph{Motivazione progettuale.} La risoluzione dell'agenzia usa un
confronto a \emph{token bidirezionale}: il marchio generico (``Autolinee
Toscane'') è un sottoinsieme dei nomi delle sotto-reti specifiche, mentre una
frase utente prolissa può esserne un soprainsieme. Inoltre, in km4city il
marchio è suddiviso in circa 40 sotto-reti; per un advisor centrato su Firenze
si preferisce la rete urbana fiorentina (\texttt{\_Agency\_888-48}), che è anche
il valore predefinito quando l'utente non nomina alcuna agenzia. Le forme dei
payload del trasporto pubblico non sono documentate con precisione: ogni parser
è quindi difensivo e registra le ``teste'' grezze dei payload in un file di log
per la calibrazione alla prima esecuzione reale.

% =============================================================================
\section{Livello MCP (\texttt{mcp\_tools.py})}
% =============================================================================
Questo modulo non implementa strumenti: (i) si connette al server e (ii) esegue
le chiamate del grafo, smussando le note stranezze di km4city. La funzione
\texttt{exec\_tool} è l'unico punto di esecuzione e \emph{non solleva mai
eccezioni}: ogni fallimento ritorna come \texttt{\{"error": \dots\}}, così il
grafo può recuperare con grazia.

\paragraph{Connessione.} \texttt{\_build\_config} scarica \texttt{/apps.json}
dalla dashboard, mantiene solo il server \texttt{snap4agentic\_advisor\_native}
e riscrive l'IP interno con \texttt{DASHBOARD\_URL}. Con un solo server nella
configurazione, FastMCP non antepone alcun prefisso ai nomi degli strumenti,
quindi si invocano con il nome ``nudo'' (es.\ \texttt{routing}).

\paragraph{Routing.} La chiamata passa per \texttt{routing\_with\_retry}, che
gestisce: (A) il corpo vuoto da cold-start con ritentativi temporizzati; (B) la
verifica dell'\texttt{error\_code} dell'envelope km4city (l'unico segnale
autorevole di successo, dato che \texttt{error\_message} è non vuoto anche in
caso di successo); (C) il rilevamento della lista di percorsi vuota (auto in zona
pedonale, origine uguale a destinazione, ecc., senza alcun codice 4xx).

\paragraph{Geocodifica.} Il geocoder di km4city non è più limitato alla regione
(indicizza anche Valencia e la Francia meridionale), perciò una ricerca
fiorentina può restituire 100 risultati esteri e zero toscani. Si applica quindi
un filtro client su un \emph{bounding box} della Toscana, una strategia a due
passaggi (prima indirizzi, poi POI come ripiego) e una ``scala'' per città
(città nominata dall'utente $>$ Firenze predefinita $>$ tutta la Toscana). Un
ulteriore ritentativo limitato copre la finestra transitoria in cui il backend
restituisce zero risultati toscani.

\paragraph{Motivazione progettuale: l'``altitudine'' della logica.} La
distinzione fra ``zona a traffico limitato'' e ``problema lato server'' non è
affidata al prompt: è calcolata in Python (funzione \texttt{\_routing\_hint})
che, in base a \texttt{routetype} e alla stringa d'errore deterministica,
produce una chiave di suggerimento. Il prompt si limita a seguirla. La logica
deterministica appartiene al codice; il prompt esprime soltanto ``leggi il campo
strutturato e agisci di conseguenza''.

% =============================================================================
\section{Client Llama4 (\texttt{llm.py}) e gestione del token}
% =============================================================================
Il client incapsula l'autenticazione e l'inferenza verso l'endpoint
\texttt{llama4-agentic-inference} (compatibile OpenAI, basato su vLLM). La
risposta è sempre un oggetto OpenAI: si ispeziona con \texttt{assistant\_message}
e \texttt{tool\_calls}. Per impostazione predefinita \texttt{tool\_choice="none"}
forza il formato OpenAI anche senza strumenti. I fallimenti transitori del
gateway (timeout upstream, errori 5xx, anche annidati in una risposta HTTP 200)
sono ritentati con backoff lineare; gli errori ``duri'' (credenziali errate,
regola API inattiva) sono sollevati subito. La chiamata bloccante è eseguita in
un thread separato (\texttt{achat}) per non bloccare l'event loop dei nodi
asincroni. Le credenziali sono lette da un file \texttt{user\_credentials.json}
(escluso dal versionamento), mai da variabili d'ambiente. Un \texttt{TokenManager}
gestisce il token OAuth (grant password / refresh), lo memorizza su disco e lo
rinnova 60 secondi prima della scadenza; sul percorso caldo (token valido in
cache) non emette alcun log, per non inondare di rumore una sessione di chat.

% =============================================================================
\section{Nodo \texttt{respond} --- formulazione testuale (LLM)}
% =============================================================================
Il terzo nodo estrae i dati del widget dall'audit, costruisce una vista
compatta dei risultati per l'LLM e produce la risposta. L'output finale segue lo
schema OpenAI: la risposta vive in \texttt{messages[-1].content}, senza alcun
campo proprietario aggiuntivo (Listato~\ref{lst:final}).

\begin{lstlisting}[caption={Forma del widget JSON restituito.}, label={lst:final}]
{
  "status": "success",          # esito in stile JSend
  "request_type": "route",      # l'intento servito
  "data": { "wkt": "...", "distance_km": 0.68, "eta": "...", ... },
  "messages": [ ... ]            # storia aggiornata per il multi-turno
}
\end{lstlisting}

\paragraph{Motivazione progettuale.}
\begin{itemize}
  \item \textbf{\texttt{temperature=0.2}} (non 0, non 0.7): a 0.7 il modello, in
  assenza di dati nei risultati, ``inventava'' numeri di linea e operatori
  inesistenti e talvolta passava all'inglese. La fedeltà ai dati viene prima
  della naturalezza: si abbassa quindi a 0.2, lasciando un minimo margine
  espressivo senza ``deragliare''.
  \item \textbf{\texttt{tool\_choice="none"} senza strumenti}: il nodo di
  formulazione non deve poter chiamare strumenti; l'opzione lo impedisce a
  livello fisico.
  \item \textbf{vista compatta dei risultati}: i payload completi (inclusa la
  geometria WKT) restano nell'audit, ma all'LLM si fornisce solo un riassunto;
  un contesto troppo grande fa degradare (o andare in errore) il backend del
  modello. In particolare le coordinate grezze sono deliberatamente nascoste al
  modello, perché non possa usarle per stimare di propria iniziativa distanze o
  tempi.
  \item \textbf{nessun campo \texttt{answer} proprietario}: l'output rispetta lo
  standard OpenAI, leggendo la risposta da \texttt{messages[-1].content}.
\end{itemize}
In caso di errore dell'LLM, una risposta di ripiego deterministica in italiano
viene comunque prodotta.

% =============================================================================
\section{Interfaccia di test (\texttt{chat.py})}
% =============================================================================
Una REPL da terminale consente il test multi-turno: si digita una domanda, si
ottiene la risposta, si prosegue (le domande di follow-up come ``\textit{e in
bus?}'' sono risolte rispetto alla storia). Vengono mostrate solo le parole
dell'LLM (\texttt{messages[-1].content}), mentre l'intero JSON del widget viene
accodato a \texttt{outputs.txt} per l'ispezione offline; i log di diagnostica
vanno in \texttt{debug.log}.

\paragraph{Motivazione progettuale.} Si è scelta una REPL da terminale e non
un'interfaccia web perché l'unico ambiente di esecuzione è JupyterHub, dove
esporre una porta web richiede un proxy fragile e il riavvio del server (con il
rischio concreto di rendere inutilizzabile l'intero singleuser server). Una REPL
da terminale, invece, funziona senza alcuna infrastruttura aggiuntiva.

% =============================================================================
\section{Limitazioni note lato server}
% =============================================================================
La sperimentazione (sonde dirette che bypassano l'orchestratore) ha permesso di
attribuire con certezza alcune limitazioni al server, distinguendole da
eventuali errori del client. La Tabella~\ref{tab:limiti} riassume lo stato
al 15/06/2026.

\begin{table}[h]
\centering
\begin{tabular}{p{3.2cm}p{3.0cm}p{6.5cm}}
\toprule
\textbf{Capacità} & \textbf{Stato} & \textbf{Evidenza / nota} \\
\midrule
Routing a piedi (centro) & OK & percorsi centrali 1--2 km corretti (WKT, distanza, tempo). \\
\addlinespace
Routing in auto & \textbf{OK (risolto)} & verificato ``Sesto $\rightarrow$ Scandicci'': percorso reale 19,12 km / 12 min con WKT completo. In passato restituiva vuoto; il lato server è stato corretto. Nessuna modifica necessaria al client. \\
\addlinespace
Routing trasporto pubblico & Da rivedere & non restituisce più un errore nudo, ma su tratte brevi il ``journey'' contiene una sola tratta a piedi (nessun mezzo reale); da riverificare sulle tratte lunghe. \\
\addlinespace
Routing a piedi (fuori centro / lunga distanza) & Limite di progetto & le tratte di media distanza fuori centro e quelle lunghe ($>$10 km) restituiscono vuoto: la copertura pedonale del backend è legata al centro città. \\
\addlinespace
Orari di fermata (timetable / realtime) & Vuoto lato server & \texttt{tpl\_stop\_timeline} restituisce sempre dizionari vuoti; anche chiamando \texttt{service\_info\_dev} con datetime ISO si ottiene HTTP 200 ma \texttt{timetable}/\texttt{realtime} vuoti: gli orari GTFS non risultano caricati lato server. Il client risponde onestamente ``orari non disponibili''. \\
\bottomrule
\end{tabular}
\caption{Capacità disponibili e limitazioni note (15/06/2026).}
\label{tab:limiti}
\end{table}

\paragraph{Nota metodologica.} Per ciascuna limitazione, la diagnosi è stata
condotta confrontando coordinate ``note buone'' e invocando lo strumento in modo
nudo (\texttt{call\_tool}), così da escludere errori del client (nome dello
strumento, parametri, valori di \texttt{routetype}). Le capacità del server
possono inoltre variare nel tempo (l'auto, ad esempio, è stata corretta dopo le
prime misurazioni): i risultati vanno periodicamente riverificati.

% =============================================================================
\section{Conclusioni}
% =============================================================================
Il client realizzato copre in modo affidabile: l'estrazione dell'intento,
il routing a piedi nel centro, il routing in auto, e la scoperta del trasporto
pubblico (linee, tratte, fermate), il tutto con gestione multi-turno e
robustezza verso le stranezze del backend km4city. Le scelte progettuali
principali --- grafo deterministico, LLM confinato a estrazione (forzata) e
formulazione (senza strumenti), logica deterministica in Python con
``altitudine'' corretta --- hanno eliminato intere classi di errori del
prototipo agentico iniziale. Restano da risolvere lato server il routing reale
del trasporto pubblico e il caricamento degli orari GTFS, oggetto delle prossime
verifiche con il referente.

\end{document}
